%=========================================================
\section*{Conclusión}
\label{cap:Conclusion}

El presente trabajo terminal logró cumplir satisfactoriamente con todos los objetivos específicos planteados, abordando de manera integral el problema de la suplantación de identidad durante la aplicación de los Exámenes a Título de Suficiencia (ETS) en la Escuela Superior de Cómputo (ESCOM). Se desarrolló un sistema funcional compuesto por una aplicación móvil que permite verificar la identidad de los estudiantes mediante reconocimiento facial, además de gestionar el acceso a los exámenes utilizando el código QR de la credencial del alumno para agilizar la entrada por parte del personal de seguridad. Asimismo, se implementó una página web que simula el Sistema de Administración Escolar (SAES), permitiendo la gestión y control de los exámenes ETS, así como de los alumnos, el personal de seguridad y el personal académico.

Durante el desarrollo del sistema, se exploraron y evaluaron distintas técnicas de reconocimiento facial, seleccionándose finalmente el modelo FaceNet debido a su alto desempeño en términos de precisión y confiabilidad. Uno de los principales retos enfrentados fue el ajuste del sistema de reconocimiento facial, el cual presentaba limitaciones iniciales al detectar rostros en orientación vertical. Este inconveniente fue resuelto mediante adecuaciones en la lógica de captura y preprocesamiento de imágenes.

Otro desafío importante fue el proceso de scraping de la credencial del alumno, el cual presentó conflictos relacionados con las dependencias de bibliotecas externas, específicamente al intentar ejecutar Selenium en el servidor en línea. Este inconveniente limitó la compatibilidad y estabilidad del entorno de despliegue. Para resolverlo, se realizaron modificaciones en la lógica interna del sistema, priorizando la funcionalidad sobre los aspectos visuales. Gracias a estos ajustes, fue posible mostrar adecuadamente la información esencial de la credencial sin comprometer el rendimiento general de la plataforma. Asimismo, se identificaron y solucionaron dificultades menores en la interfaz de usuario de algunos módulos, garantizando así una experiencia de usuario fluida y satisfactoria.

En términos generales, el sistema demostró ser una herramienta viable para fortalecer los mecanismos de control de acceso durante la aplicación de los ETS, mejorando la seguridad y reduciendo el riesgo de suplantación de identidad. Los resultados obtenidos son prometedores y sientan las bases para futuras mejoras, entre las cuales destacan la optimización del diseño de la interfaz, la integración con bases de datos institucionales en tiempo real y la ampliación del sistema para su implementación en otras unidades académicas del Instituto Politécnico Nacional.
